% Part: first-order-logic
% Chapter: completeness
% Section: downward-ls

\documentclass[../../../include/open-logic-section]{subfiles}

\begin{document}

\olfileid[hi]{fol}{com}{dls}
\olsection{लोवेनहाइम--स्कोलेम (L\"owenheim--Skolem) प्रमेय}

लोवेनहाइम--स्कोलेम (L\"owenheim--Skolem) प्रमेय कहता है कि यदि किसी
सिद्धांत का कोई अनंत प्रतिरूप है, तो उसका ऐसा प्रतिरूप भी है जो अधिक-से-अधिक
!!{denumerable} है। इस तथ्य का एक तात्कालिक परिणाम यह है कि प्रथम-क्रम तर्क
यह व्यक्त नहीं कर सकता कि किसी !!a{structure} का आकार !!{nonenumerable} है:
!!{sentence} अथवा !!{sentence}s का कोई भी समुच्चय, जिसे सभी
!!{nonenumerable} !!{structure}s संतुष्ट करती हैं, किसी !!{enumerable}
संरचना द्वारा भी संतुष्ट होता है।

\begin{thm} 
\ollabel{thm:downward-ls} यदि $\Gamma$ संगत है, तो उसका कोई
!!a{enumerable} प्रतिरूप है; अर्थात् वह ऐसी !!a{structure} में संतुष्टनीय है
जिसका प्रांत या तो परिमित है या !!{denumerable}।
\end{thm}

\begin{proof}
यदि $\Gamma$ संगत है, तो पूर्णता प्रमेय के प्रमाण से प्राप्त
!!{structure}~$\Struct M$ का प्रांत $\Domain{M}$, भाषा~$\Lang L$ के पदों के
समुच्चय से बड़ा नहीं है। अतः $\Struct M$ अधिक-से-अधिक !!{denumerable} है।
\end{proof}

\begin{thm}
\ollabel{noidentity-ls} यदि $\Gamma$, तादात्म्य-रहित प्रथम-क्रम तर्क की भाषा
में !!{sentence}s का संगत समुच्चय है, तो उसका कोई !!a{denumerable} प्रतिरूप
है; अर्थात् वह ऐसी !!a{structure} में संतुष्टनीय है जिसका प्रांत अनंत और
!!{enumerable} है।
\end{thm}

\begin{proof}
यदि $\Gamma$ संगत है और उसमें ऐसा कोई वाक्य नहीं है जिसमें तादात्म्य आता हो,
तो पूर्णता प्रमेय के प्रमाण से प्राप्त !!{structure}~$\Struct M$ का प्रांत
$\Domain{M}$, भाषा~$\Lang L'$ के पदों के समुच्चय के सर्वथा समान है। अतः
$\Struct{M}$ !!{denumerable} है, क्योंकि $\Trm[L']$ ऐसा है।
\end{proof}

\begin{ex}[स्कोलेम का विरोधाभास]
ज़र्मेलो--फ़्रेंकेल समुच्चय-सिद्धांत~$\Log{ZFC}$ अत्यंत शक्तिशाली है। इसमें
लगभग सभी गणितीय कथन व्यक्त किए जा सकते हैं, जिनमें समुच्चयों के आकारों से
जुड़े तथ्य भी हैं। उदाहरण के लिए, $\Log{ZFC}$ सिद्ध कर सकता है कि
वास्तविक संख्याओं का समुच्चय~$\Real$ !!{nonenumerable} है; वह कांटोर का
प्रमेय सिद्ध कर सकता है कि किसी भी समुच्चय का घात-समुच्चय स्वयं उस समुच्चय
से बड़ा है; इत्यादि। यदि $\Log{ZFC}$ संगत है, तो उसके सभी प्रतिरूप अनंत हैं,
और इसके अतिरिक्त उन सभी में ऐसे !!{element}s हैं जिनके बारे में सिद्धांत
कहता है कि वे !!{nonenumerable} हैं, जैसे वह अवयव जो~$\Log{ZFC}$ के इस
प्रमेय को सत्य बनाता है कि प्राकृतिक संख्याओं का घात-समुच्चय विद्यमान है।
लोवेनहाइम--स्कोलेम (L\"owenheim--Skolem) प्रमेय से $\Log{ZFC}$ के
!!{enumerable} प्रतिरूप भी हैं---ऐसे प्रतिरूप जिनमें ``!!{nonenumerable}''
समुच्चय हैं, किंतु जो स्वयं !!{enumerable} हैं।
\end{ex}

\end{document}
